\begin{textsample}{POS Dim 1 – vem\_ed\_32 – Score 11.00 – vem\_ed\_32/t168.txt}  \label{ex:f1_pos_035}
IVEC 2025: pesquisa e inovação em Intercâmbios Virtuais A equipe de Projetos Colaborativos Internacionais (PCIs) da Divisão de Extensão e Pesquisa (Depes) da Coordenadoria de Ensino Superior de Graduação (CGESG) do Centro Paula Souza (CPS) participou da International Virtual Exchange Conference (IVEC) 2025, entre os dias 15 e 17 de outubro.
O evento híbrido \textbf{é} considerado um dos mais importantes do mundo na área de \textbf{intercâmbios} virtuais / COIL.
Voltada a professores, gestores, pesquisadores, estudantes e líderes educacionais, a conferência IVEC trouxe programação extensa de palestras, mesas-redondas, \textbf{encontros} temáticos, workshops e apresentações orais.
A instituição anfitriã desta edição \textbf{foi} a Hellenic Mediterranean University (Creta, Grécia). “Este \textbf{encontro} global oferece uma plataforma para a \textbf{troca} de \textbf{conhecimentos} e a \textbf{compreensão} de novas \textbf{perspectivas} pedagógicas.
Para o CPS, \textbf{é} uma oportunidade de \textbf{destacar} as inovações \textbf{implementadas} na gestão de seus Projetos Colaborativos Internacionais, \textbf{consolidando} a instituição como \textbf{referência} e inspirando novas colaborações”, comenta Osvaldo Succi Junior, coordenador da Área de Apoio à Internacionalização do Ensino Superior da Depes / CGESG, \textbf{que} participou de dois \textbf{encontros} temáticos virtuais: “COIL Connect: Building a sustainable research group” - com os professores Jon Rubin e Daisy Ruedas (COIL Connect, EUA) e Simone Hackett (The Hague University of Applied Sciences, Países Baixos) * A peer-driven approach in COIL Innovation: The ‘Magic Makers’ Experience \textbf{for} COIL Coordinators to Thrive in a Dynamic and Evolving Landscape” - com as professoras Gabriela Méndez Carrera, consultora independente, e Rosi León (DePaul University, EUA) As professoras Neusa Haruka Sezaki Gritti e Regiane Camargo, responsáveis, respectivamente, pelos PCIs em inglês e espanhol conduzidos pela Depes/ CGESG nas Fatecs, apresentaram o workshop online “911: Tecnologia ao resgate frente a linguistic tensions e intercultural challenges”, com os professores Alejandro Molina (Universidad Católica Silva Henríquez, Chile) e Guadalupe Vadillo (Universidad Nacional Autónoma de México - UNAM).

% matched lemmas: compreensão, conhecimento, consolidar, destacar, encontro, implementar, intercâmbio, perspectiva, que, referência, ser, troca
\end{textsample}
